\documentclass[11pt,compress,t,notes=noshow, xcolor=table]{beamer}

\input{../../style/preamble_new}
\input{../../latex-math/basic-math.tex}
\input{../../latex-math/basic-ml.tex}

% ToC
\AtBeginSection[]{
	\begin{frame}[noframenumbering, plain]
	\frametitle{Outline}
	\setcounter{tocdepth}{1}
	\tableofcontents[currentsection]
\end{frame}
}

\title{Deep Learning for NLP II}

\begin{document}

\titlemeta{
Basics
}{
Introduction and Course Outline
}{
}{
\item Understand the scope of the course
\item Answers to all open question
\item Get an impression of the workload
}

% ------------------------------------------------------------------------------

\begin{frame}{people}

\vfill

\textbf{Lecturers and Tutors:}

	\begin{itemize}
		\item Prof. Dr. Valentin Hofmann
		\item Dr. Matthias Aßenmacher
		\item M.Sc. Lea Hirlimann (exercise sessions)
	\end{itemize}

\textbf{Time:}

	\begin{itemize}
		\item Lecture: Wednesday, 10.00 - 12.00
		\item Exercise: Friday, 10.00 - 12.00
	\end{itemize}

\vfill

\end{frame}


% ------------------------------------------------------------------------------

\begin{frame}{structure of the lecture}

\vfill

\textbf{Central building blocks of the lecture}

	\begin{itemize}
		\item Transformer Recap (1 week)
		\item Tokenization (1 week)
		\item Architecture Variants (3 weeks)
		\item Pre-Training, Fine-Tuning, Post-Training (1 week each)
		\item Reasoning Models (1 week)
		\item Evaluation and Benchmarking (1 week)
		\item Mechanistic Interpretability (1 week)
		\item Alignment, Bias, and Safety (1 weeks)
	\end{itemize}

\vfill

\end{frame}

% ------------------------------------------------------------------------------

\begin{frame}{structure of the exercise}

\vfill

\textbf{Coding Assignments:}

	\begin{itemize}
		\item Tokenization
		\item Retrieval-Augmented Generation
		\item Speculative Decoding
		\item LoRa Finetuning
		\item SFT vs. DPO
		\item Evaluation and Benchmarking
	\end{itemize}

\vfill

\end{frame}

% ------------------------------------------------------------------------------

\begin{frame}{structure of the exercise}

\vfill

\textbf{Reading Assignments:}

	\begin{itemize}
		\item \citebutton{Illustrated Transformer}{https://jalammar.github.io/illustrated-transformer/}
		\item \citebutton{Lost in the Middle}{https://aclanthology.org/2024.tacl-1.9/}
		\item \citebutton{H-Net}{https://openreview.net/pdf?id=ZbfLR9NbNF}
		\item \citebutton{RegMix}{https://openreview.net/forum?id=5BjQOUXq7i}
		\item \citebutton{LIMA}{https://openreview.net/forum?id=KBMOKmX2he}
		\item \citebutton{ROME}{https://proceedings.neurips.cc/paper_files/paper/2022/file/6f1d43d5a82a37e89b0665b33bf3a182-Paper-Conference.pdf}
		\item tbd
	\end{itemize}

\vfill

\end{frame}

% ------------------------------------------------------------------------------

\begin{frame}{prerequisites}
	
\vfill
	
\textit{Most important:} \textbf{Deep Learning for NLP I}

\vspace{1em}

\textbf{Machine learning basics:}

		\begin{itemize}
				\item A proper understanding of\\
				- linear algebra\\
				- loss functions\\
				- regularization\\
				- classification vs. regression\\
				- backpropagation and gradient descent\\
				- simple neural networks (MLPs)
				\item Great resource: \href{https://slds-lmu.github.io/i2ml/}{\textit{I2ML course}}
		\end{itemize}

\vfill

\end{frame}

% ------------------------------------------------------------------------------

\begin{frame}{Managing expections}

\vfill

\textbf{What to expect}

	\begin{itemize}
		\item Insights into advanced concepts and models of contemporary NLP
		\item Challenging exercises that deepen your understanding, i.e.\\
		- you will have to invest quite some time in trying to solve them\\
		- just checking our solutions won't get you far
		\item Critical discussions of the relevant research papers
		\item An (inter)actively taught course with motivated lecturers, i.e.\\
		- we expect you to actively participate in the lecture and ask questions\\
		- you won't get much out of this course by just enrolling and not showing up
	\end{itemize}


\vfill

\end{frame}

% ------------------------------------------------------------------------------

\begin{frame}{Managing expections}

\vfill

\textbf{What not to expect}

	\begin{itemize}
		\item An agentic loop engineering course
		\item Riding the hypetrain
		\item Only the latest state-of-the-art LLMs (you also need to know the relevant basics to understand what is going on)
		\item Low workload and easy exercises (you will only learn this thoroughly be investing some time)
	\end{itemize}

\vfill

\end{frame}

% ------------------------------------------------------------------------------

\begin{frame}{Managing expections}

\vfill

\textbf{Honesty}

	\begin{itemize}
		\item No intention to overload you with (unnecessary) work
		\item Building up proper skills requires \textbf{time investment}
		\item You won't get much out of the coding assignments if you just paste them into your coding agent 
		\item Up-front communication of what is going to be required
	\end{itemize}
\vfill

\end{frame}

% ------------------------------------------------------------------------------

\begin{frame}{Managing expections}

\vfill

\textbf{Workload}

	\begin{itemize}
		\item This is a 6 ECTS course
		\item 1 ECTS := 30h (i.e. this course = 180h)
		\item 13 weeks $\times$ 4h presence $\approx$ 52h
		\item 130h for preparation, self-study, \textbf{working on assignments}, and \textbf{actually reading papers}
		\item 130h / 14 weeks = \textbf{9h / week} (\textit{additionally to in-class time})
	\end{itemize}


\vfill

\end{frame}

% ------------------------------------------------------------------------------

\begin{frame}{your turn}

\vfill

\centering \Huge{\bf Any Questions?}

\vfill

\end{frame}

% ------------------------------------------------------------------------------

\endlecture
\end{document}